\documentclass[11pt]{article}
\usepackage[T1]{fontenc}
\usepackage{lmodern,amsmath,amssymb,mathtools,booktabs,longtable,array,microtype,graphicx}
\usepackage[a4paper,margin=25mm]{geometry}
\usepackage[hidelinks,hypertexnames=false]{hyperref}
\usepackage{xurl}
\newcounter{none}
\setlength{\emergencystretch}{3em}
\title{Exact Erd\H{o}s--Straus, Fable, and native-moment interfaces}
\author{The Clankers}
\date{20 September 2026}

\begin{document}
\maketitle
\begin{abstract}
This reader gives five coordinate-level continuations.  The Erd\H{o}s--Straus
part isolates an exact positive-real rank-three crossing and an integral
determinant dichotomy.  The boundary part constructs a perfect residue pairing,
identifies multiplication by \(2\eta^3\) as the trace-degenerating map, and
computes the nonconstant unit needed by the local ES/RH algebra isomorphism.
The analytic interfaces then prove the residue-to-native Gram multiplier and
its surviving relation defect, sharpen the interval-Gram inverse margin, and
derive a simultaneous finite-rank Bernstein certification procedure.  The
reader records the exact hypotheses and the remaining arithmetic and analytic
obligations; it proves neither universal Erd\H{o}s--Straus occupancy nor the
Riemann hypothesis.
\end{abstract}

\section*{Visual guide to the exact maps}

\begin{figure}[htbp]
\centering
\includegraphics[width=\textwidth]{figures/positive-real-receiving-divisor.pdf}
\caption{The positive-real family and its unique exact frame crossing.  The
curve is a locator; the rational isolating interval, Sturm count, rank-three
minor, and integral dichotomy are proved in Track~I, Sections~3--4.}
\end{figure}

\begin{figure}[htbp]
\centering
\includegraphics[width=\textwidth]{figures/residue-trace-factorization.pdf}
\caption{The perfect residue pairing, the exact intervening multiplication
map, the collision-local algebra, and the nonconstant residue unit.  Every
arrow and constant is proved in Track~II, Sections~3--5.}
\end{figure}

\begin{figure}[htbp]
\centering
\includegraphics[width=\textwidth]{figures/residue-native-gram-bridge.pdf}
\caption{The native orthogonal-polynomial factorization and the stronger
arbitrary-polynomial formula retaining its relation defect.  The multiplier,
support hypothesis, resolvent error, and defect are proved in Track~III.}
\end{figure}

\clearpage

\input{ES_RH_MULTI_CONTINUATION}

\end{document}
